ឆ្លើយតបទៅនឹងសំណួររបស់ប្រព័ន្ធផ្សព្វផ្សាយអំពីការតែងតាំងប្រធានគណៈកម្មការផ្សះផ្សាវិវាទ ដែនសមុទ្រកម្ពុជា-ថៃ ក្រោមអនុសញ្ញាអង្គការសហប្រជាជាតិស្តីពីនីតិសមុទ្រ អ្នកនាំពាក្យក្រសួងការបរទេស និងសហប្រតិបត្តិការអន្តរជាតិ សូមបញ្ជាក់ដូចខាងក្រោម៖

អ្នកផ្សះផ្សាឯករាជ្យទាំងបួនរូប បានស្នើសុំពេលវេលាបន្ថែមទៀត ដើម្បីពិចារណាឱ្យបានហ្មត់ចត់ថា គួរតែងតាំងឥស្សរជនរូបណាជាប្រធានគណៈកម្មការផ្សះផ្សា។ កម្ពុជា និងថៃ បានឯកភាពនឹងមតិស្នើសុំនេះ។

ពេលវេលាបន្ថែមដែលស្នើដោយអ្នកផ្សះផ្សា ដើម្បីពិចារណាលើការតែងតាំងប្រធានគណៈកម្មការ ផ្សះផ្សា គឺមានលក្ខណៈសមរម្យ និងសមស្របទាំងស្រុង ដោយយោងលើសារៈសំខាន់នៃការតែងតាំង ក៏ដូច ជាដំណើរការផ្សះផ្សាទាំងមូល។ ការស្នើសុំពេលវេលាបន្ថែមនេះនឹងពុំមានផលប៉ះពាល់ទៅលើថេរវេលា កំណត់ ឬជំហាននីតិវិធីណាមួយ ក្រោមអនុសញ្ញាអង្គការសហប្រជាជាតិស្តីពីនីតិសមុទ្រនោះទេ។

ការណ៍នេះឆ្លុះបញ្ចាំងពីនីតិវិធីដ៏លំបាក ដែលអ្នកផ្សះផ្សាកំពុងអនុវត្ត។ យើងសូមថ្លែងអំណរគុណ ចំពោះភាពហ្មត់ចត់របស់អ្នកផ្សះផ្សាទាំងអស់ និងសូមបញ្ជាក់ជាថ្មីនូវជំនឿទុកចិត្តរបស់យើងចំពោះដំណើរ ការនេះ។

កម្ពុជាត្រៀមជំរុញទៅមុខនូវការដោះស្រាយវិវាទដែនសមុទ្រជាមួយប្រទេសថៃដោយសន្តិវិធី ផ្អែកលើ ច្បាប់អន្តរជាតិ។